Here I'll talk about a standard set up in studying 2+1 dimensional gravity.
\subsection{Solutions to Einstein Equation in 2+1 Dimensions}

\subsubsection{Local Spacetime}

The Riemannian curvature of 2+1 dimensional gravity can be written
directly in terms of Ricci Tensor and the metric:
\begin{align}
  R_{\mu\nu\rho\sigma}=g_{\mu\rho}R_{\nu\sigma}+g_{\nu\sigma}R_{\mu\rho}-g_{\nu\rho}R_{\mu\sigma}-g_{\mu\sigma}R_{\nu\rho}-\frac{1}{2}(g_{\mu\rho}g_{\nu\sigma}-g_{\mu\sigma}g_{\nu\rho})R
\end{align}
Thus if we consider the vaccum Einstein equation $R_{\mu\nu}=0$, the
Riemannian curvature vanishes identically! And if we consider the
cosmological constant case $R_{\mu\nu}=\Lambda g_{\mu\nu}$, the
Riemannian curvature becomes:
\begin{align}
  R_{\mu\nu\rho\sigma}=\Lambda(g_{\mu\rho}g_{\nu\sigma}-g_{\mu\sigma}g_{\nu\rho})
\end{align}
This is a Riemann Tensor of a constant convature spacetime.
\begin{itemize}
  \item When $\Lambda=0$, the spacetime is locally flat, i.e.,
    locally isometric to Minkowski space $\mathbb{R}^{2,1}$.
  \item When $\Lambda>0$, the spacetime is locally de Sitter space $dS_3$.
  \item When $\Lambda<0$, the spacetime is locally anti-de Sitter space $AdS_3$.
\end{itemize}

\subsubsection{Global Spacetime}

However, being locally constant curvature spacetime doesn't mean that
the theory is trivial, it may include various global properties.
However, we shall ask what kinds of global topoloigy may admit a
constant curvature metric as well as a lorentzian structure(for
physical spacetime is lorentzian).

The answer is given by the following theorem:
\thm{
  Let $ M $ be a compact 3 manifold with a flat, time-orientable
  Lorentzian metric and a purely spacelike boundary. Then $ M $ has
  the topology of:
  \begin{align}
    M \cong [0,1] \times \Sigma
  \end{align}
  Here $ \Sigma $ is a closed surface homeomorphic to one of the
  boundary components of $ M $.
}
Thus the topology of a physical spacetime is fixed by the topology of
a initial spacelike slice $ \Sigma $. In the following context we
shall only focus on this kind of spacetime topology.

\subsection{Naive Phase Space Analysis (Witten)}

\subsubsection{Discussion Set Up}

How to quantize a Theory, we usually start from analyzing the phase
space of the theory and then construct a poisson braket of the phase
space. Then one can transform the poisson bracket into commutator and
get the quantum theory.
Thus to study the quantum gravity in 2+1 dimensions, we need to first
learn about the phase space of the classical 2+1 dimensional gravity.

The phase space of a classical theory is often defined as:
\defi{
  Phase Space

  The phase space is the space of all $ q, \dot{q} $ at $ t = 0 $ surface.
}
However, this definition is not covariant and we need tons of
constraints to make it work, thus in 2+1 dimensional gravity, we use
the following definition:
\defi{
  Phase Space

  The phase space is the space of all classical solutions to the
  equations of motion Modulo gauge transformations.
}
For 2+1 dimensional gravity, this means that we want to find all the
solutions to:
\begin{align}
  R_{\mu\nu} = 0
\end{align}
To do this, we might first assume a spacetime topology of the
manifold we want to consider, we focus on a manifold:
\begin{itemize}
  \item With topology $ M \cong [0,1] \times \Sigma $. where $ \Sigma
    $ is a genus $ g $ closed surface.
\end{itemize}

We now simplify the problem of analyzing the phase space of 2+1
dimensional gravity to finding the space of all solutions to the
vacuum Einstein equation $ R_{\mu\nu} = 0 $ on a manifold with
topology $ M \cong [0,1] \times \Sigma $. These solutions can be
found by quotienting the Minkowski space by a discrete subgroup of
the isometry group of Minkowski space, which is the Poincare group $
ISO(2,1) $. In Witten's paper he gives out a analysis of those manifolds.

\subsubsection{Manifolds with initial \& final singularities}

Witten explicitly constructs these class of solutions as examples, by
quotienting
the light cone of Minkowski, now we explain the procedure:

\begin{itemize}
  \item \textbf{Preparation: Quotienting $ \mathbb{H} $ to $ \Sigma_g $}

    We can construct a 2-manifold with genus $ g \geq 2 $ by
    quotienting the hyperbolic plane.
    \begin{itemize}
      \item \textbf{Step 1: } We first find the fundamental group of
        the surface $ \Sigma_g $ which is $ \pi_1(\Sigma_g) $.
      \item \textbf{Step 2: } We find a subgroup of $
        PSL(2,\mathbb{R}) $ naming $ \Gamma $
        which is isomorphic to $ \pi_1(\Sigma_g) $ and acts
        \textbf{discretely} on $ \mathbb{H} $.
      \item \textbf{Step 3: } The quotient $ \mathbb{H}/\Gamma  $
        defines a Riemann surface with genus $ g $ and a metric with
        constant negative curvature.
    \end{itemize}
\end{itemize}
\YL{[In fact I don't understand how this step works mathematically, eg
    how to construct the metric and why we don't require the group to
    be free but only properly continuous (discrete), etc. I need to learn
more about this step.]}
\rmk{
  Note that here we quotient by a discrete subgroup other than a
  properly discontinuous subgroup, because the action of a properly
  discontinuous subgroup is equivalent to being discrete in the case
  of being a subgroup of $ Aut(\mathbb{H}) = PSL(2,\mathbb{R}) $.
}

Then we apply a similar procedure to the light cone of Minkowski space to
construct a flat 3-manifold with topology $ M \cong [0,1] \times
\Sigma_g $ with initial singularity and final singularity.

To do this we first list out some useful facts about the 2+1
Minkowski spacetime:
\begin{itemize}
  \item Normally we have the metric of:
    \begin{align}
      ds^2 = -dt^2 + dx^2 + dy^2
    \end{align}
    And the light cone is defined by:
    \begin{itemize}
      \item $ X^+ $ future light cone: $ t^2 > x^2 + y^2, t > 0 $
      \item $ X^- $ past light cone: $ t^2 > x^2 + y^2, t < 0 $
    \end{itemize}
  \item The isometry group is the 3 dimensional Poincare group $
    ISO(2,1) $ which is the semidirect product of
    the Lorentz group $ SO(2,1) $ and the translation group $
    \mathbb{R}^{2,1} $.
    \begin{itemize}
      \item The proper orthochronous Lorentz group $ SO^+(2,1) $ is
        isomorphic to $ PSL(2,\mathbb{R}) $.
    \end{itemize}
\end{itemize}
\textbf{Quotienting the light cone to get a 3-manifold with initial \&
final singularities}

We first notice that the light cone can be made up by hyperbolic planes:
\begin{align}
  t^2 = x^2 + y^2 + \tau^2
\end{align}
where $ \tau $ is a parameter for the hyperbolic plane.
\begin{itemize}
  \item Each hyperbolic plane is isometric to $ \mathbb{H} $ and the
    isometry group of the hyperbolic plane is the Lorentz group $
    SO^+(2,1) $ which is a subgroup of the Poincare group $ ISO(2,1) $
\end{itemize}
Thus given a
$ \Sigma_g $ we can construct a subgroup of the lorentz group $ \Gamma $
isometric to $ \pi_1(\Sigma_g) $ and quotient the hyperbolic plane by
this subgroup to get a $ \Sigma_g $. If we do this quotient for the
whole $ X^+ $ light cone we can get a 3 manifold with topology $ M
\cong [0,1] \times \Sigma_g $ and with an initial singularity at $ \tau
= 0 $.
\begin{itemize}
  \item The resulting manifold is a solution to the vaccum Einstein
    Field Equation for it is contructed by quotienting an isometry of
    Minkowski spacetime. This quotient group preserves the metric
    thus the resulting manifold is still flat.
    \YL{[I don't understand why?? ]}
  \item The resulting manifold has a $ \Sigma_g \times [0,1] $
    topology. This is because the lorentz group act the hyperbolic
    plane to itself due to its feature of preserving Minkowski
    metric. Thus, $ \Gamma $ act independently on each hyperbolic
    plane. Moreover, the isometry group of the hyperbolic plane is
    the lorentz group and the hyperbolic plane is isometric to $
    \mathbb{H} $. Thus, quotienting the hyperbolic plane with the
    subgroup $ \Gamma $ of the lorentz group is equivalent to
    quotienting $ \mathbb{H} $ with the same subgroup, which gives us
    a $ \Sigma_g $. Thus, the resulting manifold has a $ \Sigma_g
    \times [0,1] $ topology.
\end{itemize}

Finally, we can count how many distinct manifold can be constructed
by this procedure, the answer depends on the choise of $ \Gamma $
subgroup. The final answer is that the moduli of genus $ g $ surface,
which is $ \mathbb{R}^{6g-6} $ gives us the label of distinct
manifolds we can construct by this procedure.
\YL{[I don't understand this counting, it need some knowledge of
Teichmuller space and moduli space of Riemann surface.]}

\subsubsection{General Flat Manifolds}

Now we try to construct more general flat manifolds with topology $ M
\cong [0,1] \times \Sigma_g $ by quotienting more general subspace of
the Minkowski space.
